\section*{Bab \ref{ch:b2:linalg} --- Aljabar Linear}

\begin{solution}{exo:b2:linalg:1}
Tiga bentuk di dalam dual berdimensi $3$: cukuplah kebebasannya. Sebuah
kaitan $\alpha\varphi_1 + \beta\varphi_2 + \gamma\varphi_3 = 0$ yang
dinilai di $(1,0,0), (0,1,0), (0,0,1)$ memberi $\alpha + \gamma = 0$,
$\alpha + \beta = 0$, $\beta + \gamma = 0$, sehingga $\alpha = \beta =
\gamma = 0$.

Basis pra-dual $(u_1, u_2, u_3)$: selesaikan $\varphi_i(u_j) =
\delta_{ij}$. Dengan menulis $u_j = (x, y, z)$: untuk $u_1$ syaratnya
$x + y = 1$, $y + z = 0$, $x + z = 0$, yang memberi $u_1 =
\bigl(\tfrac12, \tfrac12, -\tfrac12\bigr)$; setangkup dengan itu $u_2 =
\bigl(-\tfrac12, \tfrac12, \tfrac12\bigr)$, $u_3 = \bigl(\tfrac12,
-\tfrac12, \tfrac12\bigr)$.
\end{solution}

\begin{solution}{exo:b2:linalg:2}
Matriks pertama: satu-satunya permutasi berhasil kali tak nol mengirim
$1 \mapsto 3$, $2 \mapsto 2$, $3 \mapsto 1$ --- yakni \omterm{def:b2:structures:sn}{transposisi}
$(1\,3)$, bertanda $-1$: jadi \omterm{def:b2:linalg:det}{determinannya} $-abc$.

Matriks kedua: permutasi berhasil kali tak nol tidak mungkin mencampur
kedua bloknya (unsur yang menghubungkan keduanya bernilai $0$), jadi ia
terpecah menjadi permutasi $\{1,2\}$ dikali permutasi $\{3,4\}$, dan
tandanya adalah hasil kali kedua tanda itu: jumlahnya terfaktorkan menjadi
\[
(ad - bc)(eh - fg) .
\]
Aturan umum yang disarankan (dan memang benar, dengan bukti serupa):
\omterm{def:b2:linalg:det}{determinan} matriks diagonal blok adalah hasil kali \omterm{def:b2:linalg:det}{determinan}
blok-bloknya.
\end{solution}

\begin{solution}{exo:b2:linalg:3}
$F$ ditetapkan oleh dua persamaan bebas $\varphi_1(x,y,z,t)
= x + y - z - t = 0$ dan $\varphi_2 = x - 2y = 0$: menurut
\cref{thm:b2:linalg:annihilator} yang dibaca terbalik, $F^\circ =
\operatorname{Vect}(\varphi_1, \varphi_2)$ --- keduanya berada di $F^\circ$
menurut konstruksinya, keduanya bebas (tidak sebanding), dan $\dim F^\circ
= 4 - \dim F = 4 - 2 = 2$ sebab $\dim F = 2$ (dua persamaan bebas
di $\R^4$). Basisnya: $(\varphi_1, \varphi_2)$; dimensinya:
$2 + 2 = 4$, sesuai tuntutan teorema itu.
\end{solution}

\begin{solution}{exo:b2:linalg:4}
Bentuk $\varphi_i$ berjumlah $n + 1$ pada ruang berdimensi $(n+1)$, jadi
cukuplah kebebasannya. Jika $\sum_i \lambda_i \varphi_i = 0$, nilailah
pada polinomial Lagrange $L_j$ atas simpulnya: $\lambda_j = 0$. Basis
pra-dualnya adalah $(L_0, \dots, L_n)$, sebab $\varphi_i(L_j) =
L_j(a_i) = \delta_{ij}$.

Untuk bentuk integral dengan simpul $0, \frac12, 1$ pada $\R_2[X]$:
$\int_0^1 P = \sum_i c_i P(a_i)$ dengan $c_i = \int_0^1 L_i$. Hitung:
$L_0 = 2(X - \tfrac12)(X - 1)$, $\int_0^1 L_0 = \frac16$; $L_1 =
-4X(X-1)$, $\int_0^1 L_1 = \frac46$; $L_2 = 2X(X - \tfrac12)$,
$\int_0^1 L_2 = \frac16$. Karenanya
\[
\int_0^1 P = \frac{1}{6}\Bigl(P(0) + 4P\bigl(\tfrac12\bigr) +
P(1)\Bigr)
\quad (P \in \R_2[X]) :
\]
yaitu aturan Simpson, yang tepat pada polinomial kuadrat --- sebuah
pernyataan tentang \omterm{def:b2:linalg:dual}{basis dual}.
\end{solution}

\begin{solution}{exo:b2:linalg:5}
$\operatorname{im} u = Ka$ untuk suatu $a \neq 0$; maka $u(x) =
\varphi(x)\,a$ dengan $\varphi(x)$ adalah koordinat $u(x)$ pada $a$
--- yang linear pada $x$. Trace: lengkapi $a = e_1$ menjadi basis; matriks
$u$ berkolom $\varphi(e_j)\,e_1$, jadi satu-satunya unsur diagonalnya
adalah $\varphi(e_1) = \varphi(a)$, sehingga $\operatorname{tr} u =
\varphi(a)$. Lalu
\[
u^2(x) = \varphi(x)\, u(a) = \varphi(x)\varphi(a)\, a
= (\operatorname{tr} u)\, u(x).
\]
\omterm{def:b2:linalg:det}{Determinannya}, dalam dua kasus. \emph{Jika $\varphi(a) \neq 0$:} ambil
sembarang basis hiperbidang $\ker\varphi$ lalu tambahkan $a$. Maka $u$
menolkan $\ker\varphi$ (di sana $u(x) = \varphi(x)a = 0$) dan $u(a) =
\varphi(a)\,a$: matriks $I + u$ menjadi diagonal, $(1, \dots, 1,\,
1 + \varphi(a))$, sehingga $\det(I + u) = 1 + \varphi(a) = 1 +
\operatorname{tr} u$. \emph{Jika $\varphi(a) = 0$:} maka $a \in
\ker\varphi$; ambil basis $\ker\varphi$ yang vektor pertamanya
$a$, lalu tambahkan vektor $b$ dengan $\varphi(b) = 1$. Maka $I + u$
membiarkan basis $\ker\varphi$ tetap dan mengirim $b \mapsto b + a$:
segitiga dengan diagonal satu, jadi $\det(I + u) = 1 = 1 +
\operatorname{tr} u$. Kedua kasus itu cocok dengan rumusnya.
\end{solution}

\begin{solution}{exo:b2:linalg:6}
Menurut \cref{exo:b2:linalg:9}, hiperbidangnya adalah $H_A = \{M :
\operatorname{tr}(AM) = 0\}$ dengan $A \neq 0$.

\emph{Jika $A = \lambda I$:} maka $H_A$ adalah hiperbidang bertrace nol;
matriks permutasi \omterm{def:b2:structures:sn}{siklus} panjang $n$ (bernilai satu pada kedudukan $(i,
i+1)$ dan $(n, 1)$) punya invers (\omterm{def:b2:linalg:det}{determinannya} $\pm 1$ menurut
perhitungan \cref{ex:b2:linalg:permexample}) dan bertrace nol.

\emph{Jika $A$ bukan skalar:} mula-mula carilah $P$ berinvers sedemikian
sehingga $B = P^{-1}AP$ punya unsur tak nol di luar diagonal, yaitu $b_{ji}$
dengan ($j \neq i$). Memang, jika $A$ sudah punya unsur semacam itu, ambil $P = I$;
jika $A$ diagonal dengan dua unsur berbeda $d_1 \neq d_2$, mengonjugasikannya
dengan transveksi $P = I + E_{12}$ akan memunculkan unsur luar diagonal
$d_1 - d_2 \neq 0$ (hitunglah: $P^{-1}AP = A + (d_1 - d_2)E_{12}$);
sedangkan matriks diagonal yang semua unsurnya sama adalah
skalar, dan itu sudah disingkirkan. Kini ambil $M' = I + tE_{ij}$ dengan
$t = -\operatorname{tr}(B)/b_{ji}$: maka
\[
\operatorname{tr}(BM') = \operatorname{tr} B + t\,b_{ji} = 0,
\]
dan $M'$ punya invers (segitiga dengan diagonal satu). Setelah
konjugasinya dibatalkan, $M = PM'P^{-1}$ punya invers dan
$\operatorname{tr}(AM) = \operatorname{tr}(BM') = 0$: jadi $M \in H_A$.
\end{solution}

\begin{solution}{exo:b2:linalg:7}
Menurut rumus permutasi, $\det(I + tA)$ adalah polinomial dalam $t$;
suku tetapnya $1$ (ambil $t = 0$). Koefisien $t$-nya: uraikan
$\det$ sebagai \omterm{def:b2:linalg:alternating}{bentuk alternating} atas kolom $e_j + t\,c_j(A)$; menurut
kemultilinearan, suku yang linear pada $t$ mengganti tepat satu $e_j$
dengan $c_j(A)$:
\[
\sum_{j} \det(e_1, \dots, c_j(A), \dots, e_n)
= \sum_j a_{jj} = \operatorname{tr} A ,
\]
(\omterm{def:b2:linalg:det}{determinan} dengan semua kolom kanonik kecuali $c_j(A)$ pada slot
$j$ memungut unsur diagonal ke-$j$). Jadi turunannya di $0$ adalah
$\operatorname{tr} A$.

Misalkan $g(t) = \det(\eu^{tA})$. Sifat grupnya memberi $g(s + t) =
g(s)g(t)$ (kemultiplikatifan $\det$), $g$ dapat diturunkan, dan
$g'(0) = \operatorname{tr} A$ menurut uraian di atas ($\eu^{tA} = I + tA +
O(t^2)$). Morfisma $(\R, +) \to (\R^*, \times)$ yang dapat diturunkan
memenuhi $g' = g'(0)\,g$ (turunkan $g(s+t)$ terhadap $s$ di $0$), jadi
$g(t) = \eu^{t\operatorname{tr} A}$ berkat ketunggalan penyelesaian
$y' = cy$ dengan $y(0) = 1$ (jilid Tahun ke-1).
\end{solution}

\begin{solution}{exo:b2:linalg:8}
Misalkan $v_k = (1, j^k, j^{2k})^{\mathsf T}$ untuk $k = 0, 1, 2$. Dengan
memakai $1 + j + j^2 = 0$ dan $j^3 = 1$:
\[
C v_k =
\begin{pmatrix}
a + b j^k + c j^{2k}\\
c + a j^k + b j^{2k}\\
b + c j^k + a j^{2k}
\end{pmatrix}
= (a + b j^k + c j^{2k})
\begin{pmatrix} 1\\ j^k\\ j^{2k}\end{pmatrix},
\]
(periksa baris kedua: $j^k(a + bj^k + cj^{2k}) = aj^k + bj^{2k} +
cj^{3k} = c + aj^k + bj^{2k}$). Jadi $v_k$ vektor eigen dengan
nilai eigen $\lambda_k = a + bj^k + cj^{2k}$. Vektor $v_k$ membentuk basis
(Vandermonde atas $1, j, j^2$ yang berbeda), jadi $C$ dapat didiagonalkan
dengan nilai eigen tersebut, dan
\[
\det C = \lambda_0\lambda_1\lambda_2
= (a+b+c)(a + bj + cj^2)(a + bj^2 + cj).
\]
\end{solution}

\begin{solution}{exo:b2:linalg:9}
Pemetaan $\Theta \colon A \mapsto \operatorname{tr}(A\,\cdot)$ bersifat
linear dari $\mathcal{M}_n(K)$ ke dualnya, di antara dua ruang berdimensi
sama $n^2$: jadi cukuplah keinjektifannya. Jika $\operatorname{tr}(AM) =
0$ untuk setiap $M$, ambil $M = E_{ji}$: $\operatorname{tr}(A E_{ji}) =
a_{ij} = 0$ untuk setiap $i, j$, jadi $A = 0$. Maka $\Theta$ isomorfisma.

Ketunggalan trace (\cref{prop:b2:linalg:trace}): bentuk $t$ yang
menolkan semua komutator berbentuk $\operatorname{tr}(A\,\cdot)$ untuk
suatu $A$ dengan $\operatorname{tr}(A(MN - NM)) = 0$ untuk setiap $M, N$,
yakni $\operatorname{tr}((AM - MA)N) = 0$ untuk setiap $N$ (sifat
siklis), yakni $AM = MA$ untuk setiap $M$ (keinjektifan $\Theta$): jadi
$A$ komutatif dengan segalanya, sehingga ia skalar ($A$ yang komutatif
dengan semua $E_{ij}$ memaksa unsur luar diagonalnya $0$ dan unsur
diagonalnya sama), sehingga $t = c \operatorname{tr}$.
\end{solution}

\begin{solution}{exo:b2:linalg:10}
Bekerjalah atas $\C$ (matriks real nilpoten bila dan hanya bila ia
nilpoten sebagai matriks kompleks: kenilpotenan berarti $u^n = 0$).

\emph{Langkah 1: $\operatorname{tr}(u^k) = 0$ untuk $k \geq 1$.} Secara
induktif, $u^k v - v u^k = k\, u^k$: untuk $k = 1$ inilah
hipotesisnya; sedangkan untuk langkah induksinya,
\[
u^{k+1}v - vu^{k+1} = u^k(uv - vu) + (u^k v - v u^k)u
= u^{k+1} + k\,u^{k+1} .
\]
Setelah tracenya diambil: $0 = \operatorname{tr}(u^k v) -
\operatorname{tr}(vu^k) = k \operatorname{tr}(u^k)$, sehingga
$\operatorname{tr}(u^k) = 0$.

\emph{Langkah 2: trace pangkat yang nol mengakibatkan kenilpotenan (atas
$\C$).} Misalkan $\lambda_1, \dots, \lambda_r$ nilai eigen tak nol $u$
yang berbeda, dengan multiplisitas $m_1, \dots, m_r$ (pada polinomial
karakteristik, yang terurai lengkap atas $\C$ ---
\cref{ch:b2:reduction}). Trace pangkatnya adalah $\operatorname{tr}(u^k) =
\sum_i m_i \lambda_i^k$ (trigonalkan: diagonal pangkat ke-$k$ sebuah
matriks segitiga berisi pangkat ke-$k$ diagonalnya). Sistem $\sum_i m_i
\lambda_i^k = 0$ untuk $k = 1, \dots, r$ dapat dibalik secara Vandermonde
pada peubah $m_i\lambda_i$ (matriks $(\lambda_i^{k-1})$ dikali
diagonal $\lambda_i$, dengan semua $\lambda_i \neq 0$ berbeda): jadi
setiap $m_i \lambda_i = 0$, dan itu mustahil bila $m_i \geq 1$ kecuali
$r = 0$. Jadi $u$ tidak punya nilai eigen tak nol: polinomial
karakteristiknya $(-X)^n$, dan Cayley--Hamilton (\cref{ch:b2:reduction})
memberi $u^n = 0$, yakni nilpoten.
\end{solution}

\begin{solution}{exo:b2:linalg:11}
Tulis $V(a_0, \dots, a_n)$ untuk \omterm{def:b2:linalg:det}{determinannya} lalu berinduksilah pada
$n$; $V(a_0) = 1$ menjadi awalnya. Tetapkan $a_0, \dots, a_{n-1}$ lalu
pandang $D(T) = V(a_0, \dots, a_{n-1}, T)$, yakni \omterm{def:b2:linalg:det}{determinan} yang kolom
terakhirnya $(1, T, \dots, T^n)$: dengan menguraikannya sepanjang kolom
itu, $D$ adalah polinomial berderajat $\leq n$ dalam $T$ yang koefisien
$T^n$-nya adalah minor $V(a_0, \dots, a_{n-1})$. Andaikan dulu $a_0,
\dots, a_{n-1}$ berbeda. Untuk tiap $T = a_i$ ($i < n$) ada dua kolom
yang berimpit, jadi $D(a_i) = 0$: dengan $n$ akar yang berbeda dan
derajat $\leq n$,
\[
D(T) = V(a_0, \dots, a_{n-1}) \prod_{i=0}^{n-1}(T - a_i),
\]
lalu $T = a_n$ beserta hipotesis induksinya memberi rumus hasil kali
itu. Jika dua di antara $a_0, \dots, a_{n-1}$ berimpit, kedua ruasnya
bernilai $0$ (kolom yang berulang; faktor yang berulang), sehingga
rumusnya berlaku secara sepele.
\end{solution}

\begin{solution}{exo:b2:linalg:12}
\emph{$D$ punya invers.} Kalikan dari kanan dengan matriks blok
$T = \begin{psmallmatrix} I & 0\\ -D^{-1}C & I\end{psmallmatrix}$,
yang segitiga blok dengan diagonal satu, sehingga $\det T = 1$
(\omterm{def:b2:linalg:det}{determinannya}, menurut rumus permutasi, hanya memungut blok
diagonalnya --- inilah aturan blok pada \cref{exo:b2:linalg:2}):
\[
\begin{pmatrix} A & B\\ C & D\end{pmatrix} T
= \begin{pmatrix} A - BD^{-1}C & B\\ C - DD^{-1}C & D\end{pmatrix}
= \begin{pmatrix} A - BD^{-1}C & B\\ 0 & D\end{pmatrix},
\]
yang \omterm{def:b2:linalg:det}{determinannya} $\det(A - BD^{-1}C)\det D = \det\bigl((A -
BD^{-1}C)D\bigr) = \det(AD - BD^{-1}CD)$. Karena $CD = DC$, berlaku
$BD^{-1}CD = BC$: jadi \omterm{def:b2:linalg:det}{determinannya} $\det(AD - BC)$.

\emph{Untuk $D$ sembarang.} Misalkan $D_t = D + tI$; maka tetap berlaku
$CD_t = D_tC$. Keduanya, yakni
\[
f(t) = \det\begin{pmatrix} A & B\\ C & D_t\end{pmatrix}
\qquad\text{dan}\qquad
g(t) = \det(AD_t - BC)
\]
merupakan fungsi polinomial dalam $t$. Polinomial $\det(D + tI)$
bersifat monik berderajat $n$, jadi ia punya paling banyak $n$ akar:
untuk semua $t$ kecuali berhingga banyak, $D_t$ punya invers dan $f(t) =
g(t)$ menurut kasus pertama. Dua polinomial atas lapangan tak hingga
yang berimpit pada tak hingga banyak titik pastilah sama: jadi $f = g$,
dan $t = 0$ merampungkannya.
\end{solution}

\begin{solution}{pb:b2:linalg:1}
\textbf{1.} $\Phi$ bersifat linear dengan $\ker\Phi = \bigcap_i
\ker\varphi_i$ (sebuah tupel-$p$ bernilai nol bila dan hanya bila tiap
unsurnya nol). Untuk bentuk koordinat $\varepsilon_i$ pada $K^p$:
$\Phi^{\mathsf T}(\varepsilon_i) = \varepsilon_i \circ \Phi = \varphi_i$,
jadi $\operatorname{im}\Phi^{\mathsf T} \supseteq
\operatorname{Vect}(\varphi_i)$; sebaliknya
$\operatorname{im}\Phi^{\mathsf T}$ direntang oleh
$\Phi^{\mathsf T}(\varepsilon_i)$ (sebab $\varepsilon_i$ merentang
$(K^p)^*$). Jadi $\operatorname{rk}\Phi = \operatorname{rk}
\Phi^{\mathsf T} = \dim\operatorname{Vect}(\varphi_1, \dots,
\varphi_p) =: r$ (\cref{prop:b2:linalg:transposerank}), dan
rank--nulitas memberi $\dim\bigcap_i\ker\varphi_i = n - r$.

\textbf{2.} ($\Leftarrow$) Simpan satu keluarga bagian bebas yang
maksimal, katakanlah $\varphi_1, \dots, \varphi_r$, yang merentang ruang
yang sama (jadi hipotesisnya tetap berbunyi $\bigcap_{i \leq
r}\ker\varphi_i \subseteq \ker\varphi$: irisan atas semua $i$ sama dengan
irisan atas $i \leq r$, sebab tiap bentuk yang dibuang adalah sebuah
kombinasi). Pemetaan $\Psi = (\varphi_1, \dots, \varphi_r) \colon E \to
K^r$ bersifat surjektif (pertanyaan 1: ranknya $r$). Jika $\Psi(x) =
\Psi(y)$ maka $x - y \in \ker\Psi \subseteq \ker\varphi$, jadi
$\varphi(x) = \varphi(y)$: karenanya $\varphi$ terfaktorkan sebagai
$\varphi = \lambda \circ \Psi$ dengan $\lambda \colon K^r \to K$ yang
terdefinisi dengan baik; $\lambda$ bersifat linear karena $\Psi$ linear
dan surjektif (untuk $t = \Psi(x)$, $t' = \Psi(x')$: $\lambda(t + \alpha
t') = \varphi(x + \alpha x') = \lambda(t) + \alpha\lambda(t')$). Dengan
menulis $\lambda = \sum c_i \varepsilon_i$ diperoleh $\varphi = \sum_{i
\leq r} c_i\varphi_i$. ($\Rightarrow$) Jika $\varphi = \sum c_i
\varphi_i$, maka setiap $x$ yang menolkan tiap $\varphi_i$ menolkan $\varphi$.

\textbf{3.} Menurut pertanyaan 1, $\dim\bigcap\ker\varphi_i = n - r$
dengan $r = \dim\operatorname{Vect}(\varphi_i) \leq p$, dan $r = p$
bila dan hanya bila keluarganya bebas. Untuk subruang $F$ berkodimensi
$p$: \omterm{def:b2:linalg:annihilator}{anihilatornya} berdimensi $p$
(\cref{thm:b2:linalg:annihilator}); sebuah basis $(\varphi_1, \dots,
\varphi_p)$ untuk $F^\circ$ memberi $F = \bigcap_i\ker\varphi_i$ (lewat
rumus pemulihannya). Kurang dari itu tidak mungkin: irisan $q$
hiperbidang berdimensi $\geq n - q > n - p$ menurut pertanyaan 1.

\textbf{4.} Hitung $\ker\varphi_1 \cap \ker\varphi_2$: dari $x
+ y - z = 0$ dan $y + z - t = 0$, parameterkan lewat $(y, z)$: $x = z
- y$, $t = y + z$, yang memberi bidang berisi vektor $(z - y,\; y,\;
z,\; y + z)$. Di bidang itu, $\psi = x + 2y - t = (z - y) + 2y - (y + z)
= 0$: jadi menurut lema pemfaktoran $\psi \in
\operatorname{Vect}(\varphi_1, \varphi_2)$ --- dan memang $\psi =
\varphi_1 + \varphi_2$. Namun $\psi' = x + y + t = (z - y) + y +
(y + z) = y + 2z$ tidak nol secara identik di sana (dengan $y = 1, z = 0$
diperoleh $1$): jadi $\psi' \notin \operatorname{Vect}(\varphi_1,
\varphi_2)$.

\textbf{5.} Tiga bentuk pada ruang berdimensi $3$: cukuplah
kebebasannya. Jika $a\psi_0 + b\psi_1 + c\psi_2 = 0$, ujilah pada $1, X,
X^2$: $a + b + c = 0$, $b + \frac c2 = 0$, $b + \frac c3 = 0$;
mengurangkan dua yang terakhir memberi $c = 0$, lalu $b = 0$, $a = 0$.
Basis antedual: dengan menulis $P = \alpha + \beta X + \gamma X^2$ lalu
menyelesaikan $\psi_i(P_j) = \delta_{ij}$ ($P(0) = \alpha$, $P(1) =
\alpha + \beta + \gamma$, $\int_0^1 P = \alpha + \frac\beta2 +
\frac\gamma3$), diperoleh
\[
P_0 = 1 - 4X + 3X^2, \qquad
P_1 = -2X + 3X^2, \qquad
P_2 = 6X - 6X^2 .
\]
(Periksa, misalnya: $\int_0^1 P_2 = 3 - 2 = 1$, $P_2(0) = P_2(1) = 0$.)
Masalah interpolasinya terselesaikan lewat koordinat pada
basis antedual itu:
\[
P = 1\cdot P_0 + 2\cdot P_1 + \tfrac32\, P_2 = 1 + X
\]
(koefisien $X$: $-4 - 4 + 9 = 1$; koefisien $X^2$: $3 + 6 - 9 =
0$); dan memang $P(0) = 1$, $P(1) = 2$, $\int_0^1 P = \frac32$.

\textbf{6.} Kelinearan: untuk setiap $\varphi$, $J(x + \alpha
y)(\varphi) = \varphi(x + \alpha y) = J(x)(\varphi) + \alpha
J(y)(\varphi)$, yakni $J(x + \alpha y) = J(x) + \alpha J(y)$.
Keinjektifan: jika $x \neq 0$, lengkapi $x = e_1$ menjadi basis; bentuk
koordinat $e_1^*$ memberi $J(x)(e_1^*) = 1 \neq 0$. Karena $\dim E^{**} =
\dim E^* = \dim E$, keinjektifan mengakibatkan kebijektifan.

\textbf{7.} Pemuatan: untuk $x \in F$ dan $\varphi \in F^\circ$ berlaku
$J(x)(\varphi) = \varphi(x) = 0$, jadi $J(F) \subseteq
F^{\circ\circ}$. Dimensinya (dengan
\cref{thm:b2:linalg:annihilator} dua kali):
\[
\dim F^{\circ\circ} = \dim E^* - \dim F^\circ
= n - (n - \dim F) = \dim F = \dim J(F),
\]
sebab $J$ injektif. Karenanya $J(F) = F^{\circ\circ}$.

\textbf{8.} Kesamaan pertama: $\varphi$ menolkan $F + G$ bila dan hanya
bila ia menolkan $F$ maupun $G$ (ia menolkan jumlah bila dan hanya bila
ia menolkan potongannya): jadi $(F+G)^\circ = F^\circ \cap G^\circ$.
Kesamaan kedua: pemuatan $F^\circ + G^\circ \subseteq (F \cap G)^\circ$
sudah jelas (tiap sukunya menolkan $F \cap G$). Dimensinya, dengan
memakai kesamaan pertama dan Grassmann:
\[
\dim(F^\circ + G^\circ) = \dim F^\circ + \dim G^\circ -
\dim(F^\circ \cap G^\circ)
= (n - \dim F) + (n - \dim G) - \bigl(n - \dim(F +
G)\bigr),
\]
dan menurut Grassmann di $E$ itu sama dengan $n - \dim(F \cap G) =
\dim(F \cap G)^\circ$: jadi keduanya sama.

\textbf{9.} Lewat lemanya: $\ker\psi \subseteq \ker\varphi$ dengan
$p = 1$ memberi $\varphi \in \operatorname{Vect}(\psi)$, dan
$\varphi \neq 0$ membuat skalarnya tak nol. Secara langsung: pilih $x_0$
dengan $\psi(x_0) \neq 0$; setiap $x$ dapat ditulis $x = \bigl(x -
\frac{\psi(x)}{\psi(x_0)}x_0\bigr) + \frac{\psi(x)}{\psi(x_0)}
x_0$ dengan suku pertamanya di $\ker\psi = \ker\varphi$; setelah
$\varphi$ diterapkan: $\varphi(x) = \frac{\varphi(x_0)}{\psi(x_0)}\psi(x)$.

\textbf{10.} Ambil \omterm{def:b2:linalg:dual}{basis dual} $(\varphi_1^*, \dots,
\varphi_n^*)$ dari $(\varphi_1, \dots, \varphi_n)$ di dalam $E^{**}$
(\cref{def:b2:linalg:dual} yang diterapkan pada $E^*$) lalu tetapkan $u_j =
J^{-1}(\varphi_j^*)$: diperoleh basis $E$ (sebab $J$ isomorfisma,
pertanyaan 6), dengan $\varphi_i(u_j) = J(u_j)(\varphi_i) =
\varphi_j^*(\varphi_i) = \delta_{ij}$. Ketunggalan: syarat
$\varphi_i(u_j) = \delta_{ij}$ menentukan $J(u_j)$ pada basis
$(\varphi_i)$, sehingga menentukan $u_j$.

\textbf{11.} Kelinearan: $(u + \alpha v)^{\mathsf T}\psi = \psi
\circ (u + \alpha v) = u^{\mathsf T}\psi + \alpha\, v^{\mathsf
T}\psi$. Keinjektifan: jika $u \neq 0$, pilih $x$ dengan $u(x) \neq 0$
dan $\psi$ dengan $\psi(u(x)) \neq 0$ (siasat bentuk koordinat pada
pertanyaan 6): maka $u^{\mathsf T}\psi \neq 0$. Ruang $\mathcal{L}(E,F)$
dan $\mathcal{L}(F^*, E^*)$ sama-sama berdimensi $\dim E \dim F$, jadi
pemetaannya bijektif. Jika $u$ punya invers, aturan pembalikan
$(vu)^{\mathsf T} = u^{\mathsf T}v^{\mathsf T}$ memberi $u^{\mathsf
T}(u^{-1})^{\mathsf T} = (u^{-1}u)^{\mathsf T} =
\mathrm{id}_{E^*}$ dan $(u^{-1})^{\mathsf T}u^{\mathsf T} =
(uu^{-1})^{\mathsf T} = \mathrm{id}_{F^*}$, sehingga $(u^{\mathsf
T})^{-1} = (u^{-1})^{\mathsf T}$.

\textbf{12.} Untuk $x \in E$ dan $\psi \in F^*$:
\[
\bigl(u^{\mathsf T\mathsf T}(J_E x)\bigr)(\psi)
= (J_E x)\bigl(u^{\mathsf T}\psi\bigr)
= (u^{\mathsf T}\psi)(x)
= \psi\bigl(u(x)\bigr)
= \bigl(J_F(u(x))\bigr)(\psi).
\]
Karena $\psi$ sembarang, berlaku $u^{\mathsf T\mathsf T} \circ J_E = J_F
\circ u$.

\textbf{13.} Menurut \cref{prop:b2:linalg:transposerank}: $\ker
u^{\mathsf T} = (\operatorname{im} u)^\circ$, jadi $u$ surjektif
$\iff \operatorname{im} u = F \iff (\operatorname{im}u)^\circ =
\{0\}$ (\cref{thm:b2:linalg:annihilator}) $\iff u^{\mathsf T}$
injektif. Dan $\operatorname{im} u^{\mathsf T} = (\ker
u)^\circ$, jadi $u$ injektif $\iff \ker u = \{0\} \iff (\ker
u)^\circ = E^*$ $\iff u^{\mathsf T}$ surjektif.

\textbf{14.} Jika $u(F) \subseteq F$ dan $\varphi \in F^\circ$, maka
$(u^{\mathsf T}\varphi)(x) = \varphi(u(x)) = 0$ untuk $x \in F$, jadi
$u^{\mathsf T}\varphi \in F^\circ$. Sebaliknya, jika $u(F)
\not\subseteq F$, pilih $x \in F$ dengan $u(x) \notin F$; menurut
rumus pemulihan pada \cref{thm:b2:linalg:annihilator} ada
$\varphi \in F^\circ$ dengan $\varphi(u(x)) \neq 0$: maka
$(u^{\mathsf T}\varphi)(x) \neq 0$ walaupun $x \in F$, sehingga
$u^{\mathsf T}\varphi \notin F^\circ$, yakni $F^\circ$ tidak stabil.

\textbf{15.} Kita punya $u^{\mathsf T} - \lambda\,\mathrm{id}_{E^*} = (u -
\lambda\,\mathrm{id}_E)^{\mathsf T}$ (\omterm{def:b2:structures:sn}{transposisi} bersifat linear dan
$\mathrm{id}^{\mathsf T} = \mathrm{id}$), jadi kernelnya adalah
$(\operatorname{im}(u - \lambda\,\mathrm{id}))^\circ$
(\cref{prop:b2:linalg:transposerank}), yang berdimensi
\[
n - \operatorname{rk}(u - \lambda\,\mathrm{id})
= \dim\ker(u - \lambda\,\mathrm{id})
\]
menurut rank--nulitas. Khususnya kernel yang satu tak nol bila dan hanya
bila yang lain tak nol: nilai eigennya sama, multiplisitas geometriknya sama.

\textbf{16.} Pemuatan: jika $b = u(x)$ dan $u^{\mathsf T}\psi =
0$, maka $\psi(b) = \psi(u(x)) = (u^{\mathsf T}\psi)(x) = 0$, jadi
$\operatorname{im} u \subseteq (\ker u^{\mathsf T})_\circ$.
Dimensinya: untuk subruang $S \subseteq F^*$ berlaku $S_\circ =
J_F^{-1}(S^\circ)$ (uraikan: $y \in S_\circ$ bila dan hanya bila setiap
$\psi \in S$ menolkan $y$, bila dan hanya bila $J_F(y) \in S^\circ$),
jadi $\dim S_\circ = \dim F - \dim S$. Dengan $S = \ker u^{\mathsf T}$:
\[
\dim(\ker u^{\mathsf T})_\circ
= \dim F - \dim\ker u^{\mathsf T}
= \operatorname{rk} u^{\mathsf T} = \operatorname{rk} u :
\]
dimensinya sama, sehingga $\operatorname{im} u = (\ker
u^{\mathsf T})_\circ$. Dinyatakan ulang: $b \in \operatorname{im} u$ bila
dan hanya bila $\psi(b) = 0$ untuk setiap $\psi$ dengan $u^{\mathsf
T}\psi = 0$ --- itulah alternatif Fredholm.

\textbf{17.} Samakan $(K^m)^*$ dengan $K^m$ lewat $y \mapsto \psi_y$,
$\psi_y(v) = y^{\mathsf T}v$; maka $(u^{\mathsf T}\psi_y)(x) =
y^{\mathsf T}Ax = (A^{\mathsf T}y)^{\mathsf T}x$, jadi $u^{\mathsf
T}\psi_y = \psi_{A^{\mathsf T}y}$: transposnya memang matriks
\omterm{def:b2:linalg:transpose}{transpos}. \emph{Paling banyak satu:} jika $Ax = b$ dan
$A^{\mathsf T}y = 0$, maka $y^{\mathsf T}b = y^{\mathsf T}Ax =
(A^{\mathsf T}y)^{\mathsf T}x = 0 \neq 1$. \emph{Sedikitnya satu:}
jika (i) gagal, pertanyaan 16 menyediakan $\psi_y$ dengan $A^{\mathsf T}y
= 0$ dan $y^{\mathsf T}b \neq 0$; lalu skalakan $y$ agar bernilai $1$.

\textbf{18.} Di sini $A = \begin{psmallmatrix} 1 & 1 & 0\\ 0 & 1 & 1\\
1 & 2 & 1\end{psmallmatrix}$ (baris ketiga = pertama + kedua, jadi $A$
singular). Selesaikan $A^{\mathsf T}y = 0$: dari $y_1 + y_3 = 0$, $y_1
+ y_2 + 2y_3 = 0$, $y_2 + y_3 = 0$ diperoleh $y_1 = y_2 = -y_3$, yakni
garis yang direntang $y = (1, 1, -1)$. Menurut Fredholm: terselesaikan
bila dan hanya bila $y^{\mathsf T}b = b_1 + b_2 - b_3 = 0$, yakni $b_3 =
b_1 + b_2$ --- yang jelas merupakan syarat yang benar, sebab persamaan
ketiganya adalah jumlah dua yang pertama.

\textbf{19.} Matriks $L$ bernilai $1$ pada diagonalnya dan
$-\frac12$ pada kedudukan $(k, k\pm1)$ (modulo $n$): jadi setangkup,
sehingga $L^{\mathsf T} = L$ di bawah penyamaan pertanyaan 17.
\emph{Kernelnya:} jika $Lx = 0$ maka tiap $x_k = \frac12(x_{k-1} +
x_{k+1})$. Misalkan $k_0$ memaksimalkan $x_k$; rata-rata kedua
tetangganya, yang keduanya $\leq x_{k_0}$, sama dengan $x_{k_0}$ hanya
bila keduanya sama dengan $x_{k_0}$; dengan merambatkannya sepanjang
\omterm{def:b2:structures:sn}{siklus}, $x$ pastilah konstan. Sebaliknya konstanta memang tertolkan.
Jadi $\ker L^{\mathsf T} = \ker L = \R(1, \dots, 1)$, dan alternatif
Fredholm berbunyi: $Lx = b$ terselesaikan bila dan hanya bila
$(1,\dots,1)^{\mathsf T} b = \sum_k b_k = 0$ --- itulah syarat
keselarasan diskretnya: sebuah ``sebaran panas'' pada sebuah cincin
dapat diwujudkan oleh sebuah potensial bila dan hanya bila fluks
totalnya nol.

\textbf{20.} Jika $A$ antisetangkup dan $S$ setangkup, maka
\[
\operatorname{tr}(AS) = \operatorname{tr}\bigl((AS)^{\mathsf
T}\bigr) = \operatorname{tr}(S^{\mathsf T}A^{\mathsf T}) =
-\operatorname{tr}(SA) = -\operatorname{tr}(AS),
\]
jadi $2\operatorname{tr}(AS) = 0$ dan (karena $\operatorname{char} K \neq
2$) $\operatorname{tr}(AS) = 0$: yakni $\mathcal{A}_n \subseteq
\mathcal{S}_n^\circ$ (dengan menyamakan dualnya dengan matriks).
Dimensinya: $\dim\mathcal{S}_n^\circ = n^2 - \frac{n(n+1)}2 =
\frac{n(n-1)}2 = \dim\mathcal{A}_n$, jadi keduanya sama. Dengan
bertukar peran (perhitungan yang sama), $\mathcal{A}_n^\circ = \mathcal{S}_n$.

\textbf{21.} Berlaku $\operatorname{tr}(I_nM) = \operatorname{tr} M = 0$
untuk $M \in \mathfrak{sl}_n$: jadi garis $KI_n$ berada di
\omterm{def:b2:linalg:annihilator}{anihilatornya}, yang berdimensi $n^2 - (n^2 - 1) = 1$, sehingga keduanya
sama. Diterjemahkan lewat isomorfisma $A \mapsto
\operatorname{tr}(A\,\cdot)$: bentuk yang nol pada
$\mathfrak{sl}_n$ adalah $\operatorname{tr}(\lambda I_n\,\cdot) =
\lambda\operatorname{tr}$.

\textbf{22.} Misalkan $M \in \mathcal{M}_n(K)$. Polinomial $t
\mapsto \det(M - tI)$ tak nol dan berderajat $n$, jadi ia punya paling
banyak $n$ akar; $K$ berkarakteristik $0$, jadi ia tak hingga: pilihlah
$\lambda \neq 0$ yang bukan akar. Maka $M = (M - \lambda I) +
\lambda I$ menuliskan $M$ sebagai jumlah dua matriks berinvers.

\textbf{23.} \emph{Langkah 1:} untuk $P$ berinvers dan $X$
sembarang, terapkan keawetannya pada $M = XP$: $t(P(XP)P^{-1}) = t(XP)$,
yakni $t(PX) = t(XP)$. \emph{Langkah 2:} tetapkan $X$; kedua ruas
$t(BX) = t(XB)$ linear pada $B$ dan berimpit pada $B$ yang berinvers;
menurut pertanyaan 22 setiap $B$ adalah jumlah dua matriks berinvers,
jadi keduanya berimpit di mana-mana. \emph{Langkah 3:} $t$ menolkan
setiap komutator $XB - BX$; komutator merentang $\mathfrak{sl}_n$
(ditunjukkan pada bukti \cref{prop:b2:linalg:trace}), jadi $t$ nol pada
$\mathfrak{sl}_n$ dan pertanyaan 21 memberi $t =
c\operatorname{tr}$. (Sebaliknya setiap $c\operatorname{tr}$ memang awet
terhadap keserupaan: jadi trace adalah \emph{satu-satunya} invarian
linear bagi keserupaan.)

\textbf{24.} Jika $\operatorname{rk} u = r' \leq r$: ambil basis
$(f_1, \dots, f_{r'})$ untuk $\operatorname{im} u$ lalu tulis $u(x)
= \sum_{i=1}^{r'} \psi_i(x) f_i$; tiap koordinat $\psi_i(x)$ pada
$u(x)$ linear dalam $x$ (susunan $u$ dengan sebuah bentuk
koordinat), jadi $u$ adalah jumlah $r' \leq r$ pemetaan berank $\leq1$
(tambahkan nol bila perlu). Sebaliknya, jika $u = \sum_{i=1}^r
\psi_i(\cdot)f_i$, maka $\operatorname{im} u \subseteq
\operatorname{Vect}(f_1, \dots, f_r)$, yakni $\operatorname{rk} u \leq
r$. Sifat subaditifnya: tulis $u$ dengan $\operatorname{rk} u$ suku dan
$v$ dengan $\operatorname{rk} v$ suku; jumlahnya punya
$\operatorname{rk} u + \operatorname{rk} v$ suku, jadi
$\operatorname{rk}(u + v) \leq \operatorname{rk} u + \operatorname{rk} v$.

\textbf{25.} Kamusnya: sebuah subruang $F$ berpadanan dengan
$F^\circ$ yang dimensinya melengkapi
(\cref{thm:b2:linalg:annihilator}), lalu kembali lagi lewat bidualitas
(pertanyaan 6--7); jumlah bertukar dengan irisan (pertanyaan 8);
sebuah pemetaan $u$ berpadanan dengan $u^{\mathsf T}$ yang memenuhi
$\ker u^{\mathsf T} = (\operatorname{im}u)^\circ$,
$\operatorname{im}u^{\mathsf T} = (\ker u)^\circ$, dengan rank yang
sama, keinjektifan dan kesurjektifan yang bertukar, serta subruang
stabil dan nilai eigen yang berpadanan (pertanyaan 11--15); persamaan
$u(x) = b$ terselesaikan bila dan hanya bila $b$ ortogonal terhadap
$\ker u^{\mathsf T}$ (pertanyaan 16--19); dan pada $\mathcal{M}_n$
pasangan trace mewujudkan seluruh kamus itu secara konkret, dengan
trace sebagai satu-satunya invarian linear bagi keserupaan
(pertanyaan 20--23) dan rank sebagai panjang terkecil sebuah
penguraian atas tensor elementer (pertanyaan 24). Dalam dimensi tak
hingga pencacahan dimensinya gagal dan digantikan oleh hipotesis
\emph{ketertutupan} pada peta serta oleh kelengkapan --- pada ruang
Hilbert ini menjadi teorema representasi Riesz dan teori Fredholm atas
operator kompak, yang dibuktikan secara jujur pada jilid Tahun ke-3.

\end{solution}
